% Ergaenzung fuer Kapitel 7.3 / Abschnitt Resonanzansatz (eingefuegt in arxiv-hurwitz-double-sphere/main-de.tex)
% Monte-Carlo p-Wert (gefiltert): sage Alle_Kugeln_Scan.py  =>  p = 0.080 (80/1000)

\subsection{Algebraische Rauschbereinigung via Pfad-B-Idealfilter}\label{sec:empirie-resonanz}
Die Integration des quaternionischen Pfad-B-Filters am kombinatorischen Anker $M = 113160$ über die Hurwitz-Linksideale $I_5 = \mathcal{H}\cdot(2+i)$ und $I_7 = \mathcal{H}\cdot(2+i+j+k)$ eliminiert das stochastische Look-elsewhere-Rauschen des dichten Nullstellenfeldes signifikant \cite{manuskript}. Unter Anwendung des Mangoldt-$\delta$-Lifts auf Median-Höhe $Q_{50}$ kollabieren die fraktalen Resonanzordnungen der vier Moden gegen ein geschlossenes, rationales Symmetrieraster $p/q$:
\begin{equation}
    E_0 \to \frac{7}{1}, \quad E_1 \to \frac{7}{3}, \quad E_2 \to \frac{14}{11}, \quad E_3 \to \frac{7}{11}
\end{equation}
Die durchgängige Dominanz des Zählers $7$ spiegelt direkt die mathematische Invarianz unter der Idealnorm $N(I_7) = 7$ wider. Im Gegensatz zum ungefilterten Scan wird die vormalige, rein stochastische $1:1$-Kopplung von $E_2$ durch eine strukturelle $14/11$-Resonanz abgelöst, welche starr an das modulo 30 fixierte Phasenzentrum $m \equiv 11 \pmod{30}$ gekoppelt ist \cite{manuskript}.
Die Validierung über eine Monte-Carlo-Simulation (Ganzzahl-Test bei $q=1$) beweist die Selektivität des Filters: Der empirische p-Wert sinkt drastisch von $p \approx 0.865$ auf einen rauschkorrigierten Wert von $p \approx 0.080$. Während das unbeschränkte, kontinuierliche $p/q$-Raster stochastisch erwartbar eine hohe Trefferdichte aufweist, isoliert die Hurwitz-Integritätsbedingung diese diskreten Moden als intrinsische arithmetische Signaturen des Modells.
